% ------------------------------------------------------------------------
% ------------------------------------------------------------------------
% ------------------------------------------------------------------------
%                              Objetivos
% ------------------------------------------------------------------------
% ------------------------------------------------------------------------
% ------------------------------------------------------------------------

\chapter{OBJETIVOS} \label{chap:objetivos}

\subsection*{Objetivo general}

\begin{itemize}

    \item Desarrollar un algoritmo de clasificación basado en aprendizaje profundo para la
    detección de retinopatía diabética en imágenes del fondo del ojo.
    
\end{itemize}

\subsection*{Objetivos específicos}


\begin{enumerate}

    \item Realizar una revisión de la literatura científica sobre la detección de la
    retinopatía diabética utilizando técnicas de aprendizaje profundo en imágenes del
    fondo del ojo.

    \item Recopilar y organizar un conjunto de datos clínicos de imágenes del fondo del
    ojo para el entrenamiento y evaluación del modelo.

    \item Desarrollar un algoritmo para la clasificación de retinopatía diabética
    utilizando redes neuronales convolucionales.

    \item Evaluar y validar la efectividad del sistema de detección, comparándolo con
    otros sistemas análogos y a la vanguardia de la detección de retinopatía diabética a
    través de aprendizaje profundo.

\end{enumerate}

\subsection*{Trazabilidad de objetivos}

La Tabla~\ref{tab:trazabilidad-objetivos} indica en qué secciones del documento se
evidencia el cumplimiento de cada objetivo específico.

\begin{table}[H]
\centering
\renewcommand{\arraystretch}{1.3}
\caption{Trazabilidad de los objetivos específicos}
\label{tab:trazabilidad-objetivos}
\begin{tabularx}{\textwidth}{c X l}
\hline
\textbf{OE} & \textbf{Objetivo específico} & \textbf{Sección(es)} \\ \hline
1 & Realizar una revisión de la literatura científica sobre la detección de la
retinopatía diabética utilizando técnicas de aprendizaje profundo en imágenes del fondo
del ojo. &
Cap.~\ref{chap:fundamentos}, Cap.~\ref{chap:trabajos-previos} \\
2 & Recopilar y organizar un conjunto de datos clínicos de imágenes del fondo del ojo
para el entrenamiento y evaluación del modelo. &
Sec.~\ref{sec:conjuntos-datos}, Sec.~\ref{sec:preprocesamiento} \\
3 & Desarrollar un algoritmo para la clasificación de retinopatía diabética utilizando
redes neuronales convolucionales. &
Sec.~\ref{sec:arquitecturas-cnn}, Sec.~\ref{sec:config-experimental} \\
4 & Evaluar y validar la efectividad del sistema de detección, comparándolo con otros
sistemas análogos y a la vanguardia de la detección de retinopatía diabética a través de
aprendizaje profundo. &
Cap.~\ref{chap:resultados}, Sec.~\ref{subsec:comparacion-literatura},
Cap.~\ref{chap:discusion} \\ \hline
\end{tabularx}
\caption*{\textit{Fuente: Elaboración propia}}
\end{table}

\pagebreak


% ------------------------------------------------------------------------
% ------------------------------------------------------------------------ 